% Páginas de planificación: retirar de la versión final de entrega.
\section*{Reparto del fundamento teórico}
Estas páginas organizan el trabajo y se retirarán de la entrega final.
El fundamento completo del informe se construye con los cuatro apartados siguientes.
Cada integrante redacta el suyo con sus fuentes y citas.

\begingroup
\small
\renewcommand{\arraystretch}{1.3}
\noindent\begin{tabular}{@{}>{\raggedright\arraybackslash}p{1.6cm}>{\raggedright\arraybackslash}p{2.9cm}>{\raggedright\arraybackslash}p{10.5cm}@{}}
\toprule
Apartado & Responsable & Teoría que debe explicar \\
\midrule
2.1 & Belinda Toaquiza & Sistemas distribuidos: definición y características;
comparación con sistemas centralizados y paralelos; ocho transparencias solicitadas
por la guía; CAP, ejemplos y compromisos entre consistencia y disponibilidad. \\
2.2 & Nixon Piguave & Modelos de comunicación: sockets TCP/UDP, RPC, gRPC,
Protocol Buffers y RMI. Incluir diagramas de secuencia para explicar los intercambios
y una comparación conceptual de los mecanismos. \\
2.3 & Emanuel Pino & Sincronización: problema del ordenamiento de eventos,
reglas formales de Lamport, causalidad y relojes vectoriales. Incluir un ejemplo
matemático explicado. \\
2.4 & Humberto Herrera & Fallos y seguridad: fallos de caída, omisión,
temporización y bizantinos; detección y recuperación; autenticación con JWT y
OAuth 2.0. Relacionar estos aspectos con los límites del sistema. \\
\bottomrule
\end{tabular}
\endgroup

\subsection*{¿En qué archivo escribe cada uno?}
Cada persona edita \texttt{fundamento.tex} y \texttt{referencias.bib} dentro de su
carpeta en \texttt{docs/secciones/}: \texttt{belinda/}, \texttt{nixon/},
\texttt{emanuel/} o \texttt{herrera/}. El archivo principal une estos cuatro
fragmentos para formar la sección de fundamento teórico.

\subsection*{Dos responsabilidades que conviene distinguir}
\begin{itemize}
    \item \textbf{TCP:} Nixon explica su teoría en el apartado 2.2.
    Belinda programa el servidor y el cliente TCP y documenta esa implementación
    en el apartado 4.1.
    \item \textbf{Lamport:} Emanuel explica la teoría en el apartado 2.3 y verifica
    los registros en la sección 8. Belinda programa el reloj compartido que usan
    los nodos; Nixon lo integra en gRPC.
\end{itemize}

\subsection*{Fuentes y revisión}
Cada integrante busca y verifica fuentes de sus temas y cita las que realmente
utiliza. La selección final del equipo debe cumplir la guía y el enunciado:
al menos seis fuentes distintas previstas y al menos dos artículos indexados
recientes con DOI verificable. Comprobar autores, fechas y DOI antes de citar.
Todos deben comprender los cuatro apartados para la demostración.

\clearpage
\section*{Código y resto del informe}
El reparto de programación y análisis se complementa con la redacción asignada.
Cada responsable entrega su código o resultados y su explicación en LaTeX.

\begingroup
\small
\renewcommand{\arraystretch}{1.3}
\noindent\begin{tabular}{@{}>{\raggedright\arraybackslash}p{2.8cm}>{\raggedright\arraybackslash}p{6.0cm}>{\raggedright\arraybackslash}p{6.2cm}@{}}
\toprule
Responsable & Código o trabajo técnico & Otras partes que redacta \\
\midrule
Belinda Toaquiza & Configuración de nodos, reloj Lamport y registro compartidos;
servidor y cliente TCP, acuses, framing y errores.
& Objetivo e hipótesis (1); arquitectura general y diagrama TikZ (3.1);
implementación TCP (4.1). Coordina portada e integración del informe. \\
Nixon Piguave & Contrato \texttt{.proto}, generación de módulos Python y
servidor/cliente gRPC, reutilizando la base compartida.
& Contrato Protocol Buffers (3.2) e implementación gRPC (4.2).
Coordina resumen y abstract al finalizar. \\
Emanuel Pino & Runner, calentamiento, corridas y concurrencia; registro del
entorno, conservación de CSV, hashes y verificador de Lamport.
& Protocolo experimental (5) y verificación de Lamport (8).
Preguntas sobre fallo de un nodo y efecto del número de nodos (9.2 y 9.3). \\
Humberto Herrera & Script de estadísticas y rendimiento; generación de tablas
LaTeX, boxplot y barras con intervalos a partir de CSV reales.
& Resultados (6), borrador de discusión (7) y conclusiones cuantificadas (10).
Coordina las amenazas a la validez con aportes de todos. \\
\bottomrule
\end{tabular}
\endgroup

\subsection*{Tareas que realizan los cuatro}
\begin{itemize}
    \item Escribir su conclusión individual de al menos 150 palabras y registrar
    sus aportes con enlaces a commits.
    \item Revisar la discusión y la hipótesis con los resultados reales.
    Aportar a la aplicación de TCP/gRPC al PFC y su relación con CAP (9.4).
    \item Tomar su captura en su propio equipo, participar con rostro y voz en
    el video y registrar sus sesiones en la bitácora compartida.
    \item Revisar sus citas y declarar los usos reales de IA.
\end{itemize}

Belinda coordina la respuesta sobre ANSA (9.1), integrando lo que realmente
implemente el equipo. Nixon coordina guion y montaje del video; Emanuel organiza
el índice de capturas y datos. El PDF identificativo del SGA se prepara aparte.

\subsection*{Orden de integración}
Base y TCP de Belinda $\rightarrow$ gRPC de Nixon $\rightarrow$ experimentos de
Emanuel $\rightarrow$ análisis de Herrera. Mientras avanza esa cadena, todos
pueden redactar su teoría y consultar sus fuentes. Belinda reúne las secciones
ya escritas por cada responsable y comprueba la compilación.
